% !TEX root = ../../../main.tex
% 所属：第一部分：单变量分析学 / 第3章：实数的八大基本定理
% 节：3.8 单调有界定理
% 正文按原文件顺序保存；由 ../chapter3.tex 调用。
\section{单调有界定理}
在本节中我们对单调有界定理进行互证.
\begin{definition}
	对数列$\left\{x_n\right\}$,若$\forall n\in\Z_+$,有
	\begin{itemize}
		\item $x_n\leqslant x_{n+1}$,就称$\left\{x_n\right\}$是单调增的
		\item$x_n< x_{n+1}$,就称$\left\{x_n\right\}$是严格单调增的
		\item$x_n\geqslant x_{n+1}$,就称$\left\{x_n\right\}$是单调减的
		\item$x_n>" x_{n+1}$,就称$\left\{x_n\right\}$是严格单调减的
	\end{itemize}
\end{definition}
\begin{theorem}
	单调有界数列必收敛
\end{theorem}
\begin{proof}
	不妨设数列 $\left\{x_n\right\}$ 单调增加且有上界，根据确界存在定理，由 $\left\{x_n\right\}$ 构成的数集必有上确界 $\beta$ ，满足：
	
	（1）$\forall n \in \N^{+}: x_n \leqslant \beta ;$
	
	（2）$\forall \varepsilon>0, \exists x_{n_0}: x_{n_0}>\beta-\varepsilon$ ．
	
	取 $N=n_0, \forall n>N$ ：
	
	$$
	\beta-\varepsilon<x_{n_0} \leqslant x_n \leqslant \beta,
	$$
	
	
	因而 $\left|x_n-\beta\right|<\varepsilon$ ，于是得到
	
	$$
	\lim _{n \rightarrow \infty} x_n=\beta
	$$
	
	对下界的考量是同理的.
\end{proof}
我们在下一章介绍其应用.
\begin{proposition}
	单调有界定理$\Rightarrow$确界原理.
\end{proposition}

\begin{proof}
	设 $S$ 是非空有上界的实数集.
	
	取 $T=\{t \in \mathbb{Q}: t$ 是 $S$ 的上界 $\}$ ，
	显然 $T$ 非空.
	
	将 $T$ 中元素排成序列 $\left\{r_n\right\}$.
	令 $x_n=\min \left\{r_1, r_2, \ldots, r_n\right\}$ ，则 $\left\{x_n\right\}$ 单调递减，且 $x_n$ 是 $S$ 的上界）.
	
	又 $S$ 非空，取 $a \in S$ ，则 $a \leqslant x_n$ 对所有 $n$ 成立，所以 $\left\{x_n\right\}$ 有下界 $a$
	由单调有界定理，$\left\{x_n\right\}$ 收敛，设 $\lim\limits_{n\rightarrow\infty}x_n=\xi$ .
	
	下证 $\xi$ 是 $S$ 的上界：
	
	如若不然，存在 $s \in S$ 使 $s>\xi$ ，取 $\varepsilon=s-\xi>0$ ，则存在 $N$ ，当 $n \geqslant N$ 时，$x_n-\xi<s-\xi\Rightarrow x_n<s$ ，但 $x_n$ 是 $S$ 的上界，矛盾.
	
	下证 $\xi$ 是最小上界：
	
	如若不然，存在上界 $M^{\prime}$ 使 $M^{\prime}<\xi$ ，取有理数 $q$ 满足 $M^{\prime}<q<\xi$ .则 $q$ 也是 $S$ 的上界，所以 $q \in T$ ，因此 $q=r_k$ 对某个 $k$ .于是 $x_n \leqslant q$ 对所有 $n \geqslant k$ ，所以 $\lim\limits_{n\rightarrow\infty} x_n \leqslant q<\xi$ ，矛盾.
	
	因此 $\xi=\sup S$ .
\end{proof}
